\section{The Gamma quotient and its exact harmonic residues}
\label{sec:gamma}

The generating function in this section is classical at integer
arguments: it generates the Roman harmonic numbers, or complete
harmonic star sums, discussed by Sesma \cite[Eq.~(49)]{Sesma}.
We use its meromorphic continuation in the harmonic cutoff to produce
a convenient exact differential normalization of the digamma powers.
The all-depth ordinary integral below is proved directly; its
coefficient-series data are related to Young's established formulas,
and its first case is Coffey's integral for $\gamma_1$
\cite[Eq.~(1.19)]{Coffey}.

\begin{definition}\label{def:Cm}
For $x$ away from the nonpositive integers, define
\begin{equation}\label{eq:GammaGenerator}
 G(x,t)=\frac{\Gamma(1-t)\Gamma(x)}{\Gamma(x-t)}
       =\sum_{m\ge0}\CM_m(x)\frac{t^m}{m!}.
\end{equation}
Coefficient extraction is at $t=0$; the resulting functions of $x$
are extended meromorphically.
\end{definition}

Put
\begin{equation}\label{eq:hr}
 h_1(x)=\psi(x)+\gamma,\qquad
 h_r(x)=\zeta(r)-\zeta(r,x)
       =\zeta(r)+\frac{(-1)^{r+1}}{(r-1)!}\psi^{(r-1)}(x)
       \quad(r\ge2).
\end{equation}
The log-Gamma Taylor expansion yields
\begin{equation}\label{eq:logG}
 \log G(x,t)=\sum_{r\ge1}h_r(x)\frac{t^r}{r}.
\end{equation}
Consequently $\CM_m=Y_m(h_1,1!h_2,\ldots,(m-1)!h_m)$,
where $Y_m$ is the complete exponential Bell polynomial. In particular
\begin{align}
 \CM_0&=1,\qquad \CM_1=h_1,\qquad
 \CM_2=h_1^2+h_2,\\
 \CM_3&=h_1^3+3h_1h_2+2h_3,\\
 \CM_4&=h_1^4+6h_1^2h_2+3h_2^2+8h_1h_3+6h_4,\\
 \CM_5&=h_1^5+10h_1^3h_2+15h_1h_2^2+20h_1^2h_3
                +20h_2h_3+30h_1h_4+24h_5.
 \label{eq:Cfirst}
\end{align}
At every degree a finite recurrence suffices:
\begin{equation}\label{eq:Crecurrence}
 \CM_{m+1}=
 \sum_{r=0}^m\binom mr r!h_{r+1}\CM_{m-r}.
\end{equation}
At a positive integer $n$ the generator is
$\prod_{k=1}^{n-1}(1-t/k)^{-1}$, so $\CM_m(n)/m!$ is the
weakly ordered harmonic sum of $m$ unit indices.

\begin{theorem}[Exact poles and removable values]\label{thm:Cpoles}
For $m\ge2$, $\CM_m$ has only simple poles, with
\begin{equation}\label{eq:Cresidue}
 \Res_{x=-n}\CM_m(x)=-m!e_{m-1}(n),\qquad n\ge0.
\end{equation}
The points $0,-1,\ldots,-(m-2)$ are removable. Its first actual pole
is $-(m-1)$ and has residue $-m$. At the removable points,
\begin{equation}\label{eq:Cnegative}
 \boxed{\displaystyle
 \CM_m(-j)=-m!\sum_{k=0}^j e_k(j)\zeta(m-k),
       \qquad 0\le j\le m-2.}
\end{equation}
In particular $\CM_m(0)=-m!\zeta(m)$, and $\CM_m(1)=0$.
\end{theorem}
\begin{proof}
The reciprocal Gamma factor is entire, so the Gamma quotient has
at most simple poles in $x$. The Gamma recurrence gives exactly
\[
 \Res_{x=-n}G(x,t)
 =\frac{(-1)^n\Gamma(1-t)}{n!\Gamma(-n-t)}
 =-t\prod_{k=1}^n\left(1+\frac tk\right).
\]
Taking coefficients proves the residue assertion. Since $e_{m-1}(n)$
vanishes precisely for $n<m-1$ and is positive thereafter, it also
proves the removability and the first-pole statements.

For the values, set $x=-j+\varepsilon$ and
$E_j(t)=\prod_{k=1}^j(1+t/k)$. Expansion of both Gamma factors gives
the constant coefficient in $\varepsilon$ as
\[
 -tE_j(t)\bigl(H_j-\gamma-\psi(-j-t)\bigr).
\]
By the digamma recurrence this equals
\begin{equation}\label{eq:Cnegativeproof}
 tE_j(t)(\gamma+\psi(1-t))
 -tE_j(t)H_j+E_j(t)+tE_j(t)\sum_{k=1}^j\frac1{k+t}.
\end{equation}
All terms after the first are polynomials of degree at most $j+1$.
For the last term this follows by cancelling the corresponding factor
of $E_j$. Finally
$\gamma+\psi(1-t)=-\sum_{r\ge1}\zeta(r+1)t^r$.
At degree $m\ge j+2$ only this series contributes, proving
\eqref{eq:Cnegative}. Since the pole coefficient at that degree is
zero, its finite Laurent value is the ordinary removable value.
The identity at $1$ follows from $G(1,t)=1$.
\end{proof}

\begin{theorem}[A normally convergent partial-fraction expansion]
\label{thm:CML}
For $m\ge2$, away from its actual poles,
\begin{equation}\label{eq:CML}
 \boxed{\displaystyle
 \CM_m(z)=-m!\zeta(m)-m!\sum_{n\ge1}e_{m-1}(n)
       \left(\frac1{n+z}-\frac1n\right).}
\end{equation}
The series converges normally on every compact subset of that domain.
Terms whose coefficient $e_{m-1}(n)$ is zero are omitted, including
at a removable point where their displayed denominator vanishes.
\end{theorem}
\begin{proof}
Since $e_{m-1}(n)=O((1+\log n)^{m-1})$, the subtracted series
converges normally. To determine the entire remainder, integrate
$z\CM_m(w)/(w(w-z))$ on $|w|=N+1/2$. The right-half-plane
digamma asymptotics and their derivatives, combined with reflection
on the left, give $\CM_m(w)=O((1+\log N)^m)$ uniformly on these
circles. The distance to the nonpositive integral poles stays bounded
below. In the reflection formula the cotangent and its fixed
derivatives are bounded there, while the remaining polygammas have
their usual sectorial bounds \cite{DLMF}. The integral tends to zero
like $O((1+\log N)^m/N)$.

Its residues at $w=z$ and $w=0$ are $\CM_m(z)$ and
$-\CM_m(0)$. Each pole $w=-n$ contributes
$-m!e_{m-1}(n)z/(n(n+z))$. The residue theorem and
Theorem~\ref{thm:Cpoles} give \eqref{eq:CML}. Thus matching the
principal parts alone is not being used to discard an arbitrary
entire function.
\end{proof}

\subsection{An ordinary integral isolates one strict-depth coefficient}

\begin{theorem}[The all-depth ordinary-integral bridge]
\label{thm:Cintegral}
For every integer $m\ge2$,
\begin{equation}\label{eq:Cintegral}
 \boxed{\displaystyle
 \int_0^1\CM_m(x)\,dx=m!\bigl(c_m-\lambda_{m-1}\bigr).}
\end{equation}
The integrand extends continuously to both endpoints. There is no
finite-part operation in this identity.
\end{theorem}
\begin{proof}
Use the opposite-sign parameter
$F(x,t)=G(x,-t)=\Gamma(1+t)\Gamma(x)/\Gamma(x+t)$ and define
\[
 \Phi(t)=\int_0^1\{F(x,t)-1+t(\psi(x)+\gamma)\}\,dx.
\]
Its origin residue is $t-t=0$. More precisely, the subtracted
integrand is jointly holomorphic in $t$ near zero and continuous in
$x\in[0,1]$; this follows by subtracting $t/x$ from the Gamma
quotient. Uniform coefficient extraction on $|t|\le\rho<1$ gives
\begin{equation}\label{eq:PhiTaylor}
 \Phi(t)=\sum_{m\ge2}\frac{(-1)^m t^m}{m!}
                    \int_0^1\CM_m(x)\,dx.
\end{equation}

For $\re t>0$, Euler's beta integral and the standard digamma
integral give, after combining their integrands,
\[
 \Phi(t)=t\int_0^\infty
 \left\{\frac{(1-e^{-v})^t-1}{v}
             +\frac{e^{-v}}{1-e^{-v}}\right\}dv-1.
\]
The combination is integrable at both ends. For example the original
integrand before integration in $x$ is
$e^{-xv}((1-e^{-v})^{t-1}-(1-e^{-v})^{-1})
+e^{-v}/(1-e^{-v})$; its joint integral is absolutely convergent
when $\re t>0$, which justifies Fubini.
Integrate the second term from a lower cutoff $\delta$ and use
$\int_\delta^\infty e^{-v}/(1-e^{-v})\,dv
=-\log(1-e^{-\delta})$. Comparing with
$\int_\delta^1(v^t-1)\,dv/v$ yields
\begin{equation}\label{eq:PhiMellin}
 \Phi(t)=t\int_0^\infty
 \frac{(1-e^{-v})^t-1-\boldsymbol1_{(0,1)}(v)(v^t-1)}v\,dv.
\end{equation}
This subtracted integral is holomorphic for $\re t>-1$ locally
about zero. Near zero, $(1-e^{-v})^t-v^t=O(v^{\re t+1})$
uniformly on compact parameter sets; at infinity the numerator
decays exponentially. Thus the equality continues to $|t|<1$.

At degree $m\ge2$, its right side has coefficient
\[
 \frac{(-1)^{m-1}}{(m-1)!}
             \int_0^\infty\frac{V_{m-1}(v)}v\,dv.
\]
Comparison with \eqref{eq:PhiTaylor} gives
$\int_0^1\CM_m=-m I_{m-1,0}$.
Equation~\eqref{eq:lambda_integral} proves the claimed formula.
\end{proof}

The first two instances are
\begin{align}
 \int_0^1\{(\psi+\gamma)^2+\zeta(2)-\psi'\}\,dx
   &=\gamma^2-\zeta(2)+2\gamma_1,\label{eq:C2integral}\\
 \int_0^1\{(\psi+\gamma)^3
      +3(\psi+\gamma)(\zeta(2)-\psi')
      +2\zeta(3)+\psi''\}\,dx
   &=6(c_3-\eta).
 \label{eq:C3integral}
\end{align}
The apparent singularities in each complete integrand cancel; its
individual displayed terms should not be separately integrated as
ordinary improper integrals. The first formula rearranges exactly to
Coffey's known identity cited above.

\subsection{The known Cauchy-series coordinates recovered from the same quotient}

An additional beta representation is useful both for interpretation
and for stable numerical evaluation:
\begin{equation}\label{eq:Cbeta}
 \CM_m(x)=(x-1)\int_0^1(-\log y)^m(1-y)^{x-2}\,dy,
       \qquad m\ge1,\quad\re x>1-m.
\end{equation}
For $\re x>1$ it follows by differentiating
$G(x,t)=(x-1)B(1-t,x-1)$. The integral on the right then continues
it to the stated half-plane: near $y=1$ the exponent is $x+m-2$.
This also independently displays the location of the first pole.

Define positive rational coefficients
\begin{equation}\label{eq:an}
 a_n=\frac1{(n-1)!}\int_0^1(r)_n\,dr,
 \qquad a_1=\frac12,\quad a_2=\frac56,\quad a_3=\frac98,
\end{equation}
where $(r)_n=r(r+1)\cdots(r+n-1)$ is a rising factorial. If
$L=-\log(1-y)$, their generating kernel is
\begin{equation}\label{eq:Kkernel}
 K(y)=\sum_{n\ge1}a_ny^{n-1}
      =\frac{L-y}{(1-y)^2L^2}.
\end{equation}
Indeed integrate $(1-y)^{-r}=\sum (r)_ny^n/n!$ over $r\in[0,1]$
and differentiate the resulting $y/((1-y)L)$.
Integrating \eqref{eq:Cbeta} in $x$, for $m\ge2$, gives
\[
 \int_0^1\CM_m(x)\,dx
   =-\int_0^1(-\log y)^m K(y)\,dy.
\]
The integrals are absolutely convergent; alternatively Tonelli applies
to their nonnegative negatives. The elementary Mellin integral of
$y^{n-1}$ therefore yields
\begin{equation}\label{eq:knowncauchy}
 \boxed{\displaystyle
 \lambda_d=c_{d+1}+\sum_{n\ge1}\frac{a_n}{n^{d+2}},
       \qquad d\ge1.}
\end{equation}
This coefficient series is an equivalent form of the established
Cauchy/Bernoulli-second-kind series in Young's work
\cite{Young2023,Young2024}; it is included as an independent
normalization check and a connection with the meromorphic differential
realization, not as a new coefficient formula.
